\documentclass[11pt]{article}

% ======================
% Packages
% ======================
\usepackage[utf8]{inputenc}
\usepackage[T1]{fontenc}
\usepackage{lmodern}
\usepackage{amsmath, amssymb, amsthm}
\usepackage{geometry}
\usepackage{listings}
\usepackage{xcolor}
\usepackage{graphicx}
\usepackage{algorithm}
\usepackage{algpseudocode}
\usepackage{hyperref}
\usepackage{enumitem}
\usepackage{enumerate}


\usepackage{graphicx}    % For \includegraphics
\usepackage{tikz}        % For drawing graphs and diagrams
\usepackage{amsmath}     % For math symbols (e.g., \infty, \emptyset)
\usepackage{amssymb}     % Additional math symbols
\usepackage{array}       % For table formatting
\usepackage{booktabs}    % For better table rules (optional, but nice)
\usepackage{xcolor}      % For colored text (if needed)
\usepackage{caption}     % For \captionof{} (optional, for Beamer)


% ======================
% Page Layout
% ======================
\geometry{a4paper, margin=1in}
\setlength{\parindent}{0pt}
\setlength{\parskip}{1em}

% ======================
% Colors & Code Styling
% ======================
\definecolor{codegreen}{rgb}{0,0.6,0}
\definecolor{codegray}{rgb}{0.5,0.5,0.5}
\definecolor{codepurple}{rgb}{0.58,0,0.82}
\definecolor{backcolour}{rgb}{0.95,0.95,0.92}

\lstdefinestyle{mystyle}{
    backgroundcolor=\color{backcolour},
    commentstyle=\color{codegreen},
    keywordstyle=\color{magenta},
    numberstyle=\tiny\color{codegray},
    stringstyle=\color{codepurple},
    basicstyle=\ttfamily\small,
    breakatwhitespace=false,
    breaklines=true,
    captionpos=b,
    keepspaces=true,
    numbers=left,
    numbersep=5pt,
    showspaces=false,
    showstringspaces=false,
    showtabs=false,
    tabsize=2
}

\lstset{style=mystyle}

% ======================
% Theorem-like Environments
% ======================
\newtheorem{problem}{Problem}
\newtheorem{subproblem}{Subproblem}[problem]
\newtheorem{solution}{Solution}[problem]

% ======================
% Title Info
% ======================
\title{\textbf{Fundamental Algorithm Techniques} \\ Problem Set \#8} % \author{}
\date{\textcolor{red}{Review on April 2}} %Due: November 01, 2025}

% ======================
% Document
% ======================
\begin{document}

\maketitle



\begin{problem}[SCC and reversal, 4/10 pts]
Let $G$ a directed graph: 
\begin{enumerate}
    \item Describe an algorithm to compute the reversal $\operatorname{rev}(G)$ of a directed graph $G$ in $O(V + E)$ time.
    
    \item Prove that for every directed graph $G$, the strong component graph $\operatorname{scc}(G)$\footnote{scc(G) merges every strong component into a single point} is acyclic.
    
    \item Prove that $\operatorname{scc}(\operatorname{rev}(G)) = \operatorname{rev}(\operatorname{scc}(G))$ for every directed graph $G$.
    
    \item Fix an arbitrary directed graph $G$. For any vertex $v$ of $G$, let $S(v)$ denote the strong component of $G$ that contains $v$. For all vertices $u$ and $v$ of $G$, prove that $u$ can reach $v$ in $G$ if and only if $S(u)$ can reach $S(v)$ in $\operatorname{scc}(G)$.
\end{enumerate}
\end{problem}


\begin{problem}[Euler Tour, 2/10 pts]

An Euler tour of a strongly connected, directed graph $G = (V, E)$ is a cycle that traverses each edge of $G$ exactly once (independantly of starting point), although it may visit a vertex more than once.

\begin{enumerate}
    \item Show that $G$ has an Euler tour if and only if $\text{in-degree}(v) = \text{out-degree}(v)$\footnote{in-/out-degree number or in/out-flowing arrows respectively.} for each vertex $v \in V$.
    
    \item Describe an $O(E)$-time algorithm to find an Euler tour of $G$ if one exists. (Hint: Merge edge-disjoint cycles.)
\end{enumerate}

\end{problem}

\begin{problem}[Topological Search, 4/10 pts]
Perform Topological Sort on course example: \\
\{$A\rightarrow B, A\rightarrow C, B \rightarrow C, B \rightarrow D, C \rightarrow E,  D \rightarrow E, D \rightarrow F, G \rightarrow F, G \rightarrow E\}$.

Starting from A, find orders of the course.
Start from another node at random and find one or two sortings\footnote{\href{https://www.youtube.com/watch?v=xqLY-ZLJig8}{video topological sort}...}.



\end{problem}



\end{document}

